%
% notification-mail-event-markread.tex  (FROZEN NEGATIVE CONTROL: do not edit)
%
% ReadCommit marks the mail read when it announces it.
%
% Z Specification for the mail and wall channel events as the final switch makes
% them: each event carries the arrival's high-water mark, the mark ordering what
% was announced; the arrivals reading is read-only; and the daemon's NATS push
% for mail and wall feeds the same poll path the poller already runs.
%
% Auxiliary to notification.tex (Channel Delivery, NotifyBelt, the pendingNotify
% mutex) in the manner of notification-race.tex: notification.tex carries thirty
% components and its checks do not terminate in reasonable time, so this file lifts
% out the few components the claims depend on.  It does not extend it.
%
% Type-check: fuzz -t notification-mail-event.tex
% Model-check: see the end of this file.
%

\documentclass[a4paper,10pt,fleqn]{article}
\usepackage[margin=1in]{geometry}
\usepackage{fuzz}
\usepackage[colorlinks=true,linkcolor=blue,citecolor=blue,urlcolor=blue]{hyperref}

\begin{document}

\title{Biff Mail and Wall Channel Events: A Z Specification}
\author{Punt Labs}
\date{November 2026}
\maketitle

%%----------------------------------------------------------------------
\section{Introduction}
%%----------------------------------------------------------------------

The channel emitter announces an arrival to a Claude Code session with one
event per kind (mail, wall), carrying a count for the text.  What it compares is
not the count but the relay's \emph{high-water mark} for the kind: mail is newer
when the mark passes the last-sent one, and the record advances to the mark the
reading took before the send (\texttt{Arrivals}).  Three claims hold of it, and
the final switch relies on them:

\begin{enumerate}
  \item \textbf{At most one event per arrival.}  Each announced arrival leads to
    at most one channel event of its kind, however often the poll path runs and
    however many times the daemon pushes.
  \item \textbf{The arrivals reading is read-only.}  Reading the mark and the count
    marks no mail read; only the user's \texttt{read} does.
  \item \textbf{Push and poll are one path.}  With biffd the daemon pushes a
    notice of a mail or wall arrival to the client over its subscription, and the
    client answers by waking the poller, which runs the same read, decide and
    commit; the periodic poll runs it too.  An arrival is announced by push or by
    poll, at least once, and a duplicated or lost push changes nothing.
\end{enumerate}

The serialisation is \texttt{ChannelEmitter}'s lock: a reading is taken, decided
and committed with the lock held, so two wakes cannot both find the same arrival
newer.  The rate window is outside the model: it only delays an event, and the
property is about the events that do go out.

%%----------------------------------------------------------------------
\section{Basic Types}
%%----------------------------------------------------------------------

\begin{zed}
  ZBOOL ::= ztrue | zfalse
\end{zed}

\begin{zed}
  Kind ::= kMail | kWall
\end{zed}

\begin{zed}
  Reader ::= rPush | rPoll
\end{zed}

\begin{axdef}
  maxMark : \nat
  \where
  maxMark = 2
\end{axdef}

$Reader$ names the two ways a wake reaches the emitter, so that a reading
started by each can be in flight at once.  For a wall the mark counts the keys
that were never announced.

%%----------------------------------------------------------------------
\section{State}
%%----------------------------------------------------------------------

\begin{schema}{Relay}
  mark : Kind \fun (0 \upto maxMark) \\
  arrivedN : 0 \upto maxMark \\
  readN : 0 \upto maxMark \\
  unreadN : 0 \upto maxMark
\end{schema}

The relay: the high-water mark of each kind, which rises on an arrival and never
falls, and the mail counts.  $unreadN$ is what the relay holds unread; the model
keeps it apart from $arrivedN - readN$ so that a step that marks mail read
without a read is a state.

\begin{schema}{Emitter}
  sentMark : Kind \fun (0 \upto maxMark) \\
  wake : \power Kind \\
  locked : ZBOOL \\
  busy : Reader \fun ZBOOL \\
  rkind : Reader \fun Kind \\
  rsnap : Reader \fun (0 \upto maxMark) \\
  rnewer : Reader \fun ZBOOL \\
  rbase : Reader \fun (0 \upto maxMark) \\
  cover : (Kind \cross (1 \upto maxMark)) \fun (0 \upto 2)
\end{schema}

The emitter: the last-sent mark per kind, the kinds a wake asked it to read, the
lock, and, per reader, the reading in flight: its kind, the last-sent mark it
found, the mark it took, and whether it found the kind newer.  $cover~(k, n)$ counts the events of kind $k$
that announced arrival $n$ (a ghost).

\begin{schema}{MailEventState}
  Relay \\
  Emitter
\end{schema}

\begin{schema}{MailEventInv}
  MailEventState
  \where
  \forall k : Kind @ sentMark~k \leq mark~k \\
  \forall k : Kind; n : 1 \upto maxMark @ cover~(k, n) \leq 1 \\
  \forall k : Kind; n : 1 \upto maxMark @
    ( cover~(k, n) = 1 \iff n \leq sentMark~k ) \\
  unreadN + readN = arrivedN \\
  \forall r : Reader @ busy~r = ztrue \implies locked = ztrue
\end{schema}

The properties.  Nothing is announced beyond the relay's mark.  No arrival is
announced twice, and the arrivals announced are exactly those up to the last-sent
mark.  The mail the relay holds unread is what has arrived and was not read: an
announcement marks nothing read.  A reading in flight holds the lock.

\begin{schema}{Invariant}
  MailEventInv
\end{schema}

\begin{schema}{Init}
  MailEventState'
  \where
  mark' = \{ kMail \mapsto 0, kWall \mapsto 0 \} \\
  arrivedN' = 0 \\
  readN' = 0 \\
  unreadN' = 0 \\
  sentMark' = \{ kMail \mapsto 0, kWall \mapsto 0 \} \\
  wake' = \emptyset \\
  locked' = zfalse \\
  busy' = \{ rPush \mapsto zfalse, rPoll \mapsto zfalse \} \\
  rkind' = \{ rPush \mapsto kMail, rPoll \mapsto kMail \} \\
  rsnap' = \{ rPush \mapsto 0, rPoll \mapsto 0 \} \\
  rnewer' = \{ rPush \mapsto zfalse, rPoll \mapsto zfalse \} \\
  rbase' = \{ rPush \mapsto 0, rPoll \mapsto 0 \} \\
  cover' = \{ k : Kind; n : 1 \upto maxMark @ (k, n) \mapsto 0 \}
\end{schema}

%%----------------------------------------------------------------------
\section{Operations}
%%----------------------------------------------------------------------

\begin{schema}{Arrive}
  \Delta Relay \\
  \Xi Emitter \\
  k? : Kind
  \where
  mark~k? < maxMark \\
  mark' = mark \oplus \{ k? \mapsto mark~k? + 1 \} \\
  arrivedN' = ( \IF k? = kMail \THEN arrivedN + 1 \ELSE arrivedN ) \\
  unreadN' = ( \IF k? = kMail \THEN unreadN + 1 \ELSE unreadN ) \\
  readN' = readN
\end{schema}

A mail or a wall arrives at the relay and the mark of its kind rises.  A wall
changes no mail count.

\begin{schema}{PushWake}
  \Xi Relay \\
  \Delta Emitter \\
  k? : Kind
  \where
  mark~k? > sentMark~k? \\
  wake' = wake \cup \{ k? \} \\
  sentMark' = sentMark \\
  locked' = locked \\
  busy' = busy \\
  rkind' = rkind \\
  rsnap' = rsnap \\
  rnewer' = rnewer \\
  rbase' = rbase \\
  cover' = cover
\end{schema}

The daemon's push of an arrival notice reaches the client, which asks the poller
for a pass of that kind.  It is enabled again and again for the same arrival: a
push is at least once, and a wake is a set, so a repeat changes nothing.  A push
that is lost is the absence of this step, and the next operation covers it.

\begin{schema}{PollWake}
  \Xi Relay \\
  \Delta Emitter
  \where
  wake' = Kind \\
  sentMark' = sentMark \\
  locked' = locked \\
  busy' = busy \\
  rkind' = rkind \\
  rsnap' = rsnap \\
  rnewer' = rnewer \\
  rbase' = rbase \\
  cover' = cover
\end{schema}

The periodic wake of the poller (the five-second liveness check) asks for every
kind, with or without a push.

\begin{schema}{ReadBegin}
  \Xi Relay \\
  \Delta Emitter \\
  r? : Reader \\
  k? : Kind
  \where
  k? \in wake \\
  busy~r? = zfalse \\
  locked = zfalse \\
  wake' = wake \setminus \{ k? \} \\
  locked' = ztrue \\
  busy' = busy \oplus \{ r? \mapsto ztrue \} \\
  rkind' = rkind \oplus \{ r? \mapsto k? \} \\
  rsnap' = rsnap \oplus \{ r? \mapsto mark~k? \} \\
  rnewer' = rnewer \oplus
    \{ r? \mapsto ( \IF mark~k? > sentMark~k? \THEN ztrue \ELSE zfalse ) \} \\
  rbase' = rbase \oplus \{ r? \mapsto sentMark~k? \} \\
  sentMark' = sentMark \\
  cover' = cover
\end{schema}

The read and decide step, with the lock: the reading takes the mark, and decides
the kind newer when the mark passes the last-sent one.  It reads and changes no
mail count ($\Xi Relay$).

\begin{schema}{ReadCommit}
  \Delta Relay \\
  \Delta Emitter \\
  r? : Reader
  \where
  busy~r? = ztrue \\
  mark' = mark \\
  arrivedN' = arrivedN \\
  readN' = readN \\
  unreadN' = 0 \\
  wake' = wake \\
  locked' = zfalse \\
  busy' = busy \oplus \{ r? \mapsto zfalse \} \\
  rkind' = rkind \\
  rsnap' = rsnap \\
  rnewer' = rnewer \\
  rbase' = rbase \\
  sentMark' = sentMark \oplus
    \{ rkind~r? \mapsto ( \IF rnewer~r? = ztrue \land
         rsnap~r? > sentMark~(rkind~r?)
       \THEN rsnap~r? \ELSE sentMark~(rkind~r?) ) \} \\
  cover' = cover \oplus
    ( \IF rnewer~r? = ztrue \THEN
        \{ n : 1 \upto maxMark | rbase~r? < n \land
             n \leq rsnap~r? @
             (rkind~r?, n) \mapsto cover~(rkind~r?, n) + 1 \}
      \ELSE \emptyset )
\end{schema}

The commit: if the reading found the kind newer, one event goes out that
announces the arrivals past the last-sent mark it found, up to the mark it took,
and the record advances to that mark.  Either way the lock is released.

\begin{schema}{ClientRead}
  \Delta Relay \\
  \Xi Emitter
  \where
  unreadN > 0 \\
  mark' = mark \\
  arrivedN' = arrivedN \\
  readN' = readN + 1 \\
  unreadN' = unreadN - 1
\end{schema}

The user's own \texttt{read} marks mail read.  It is the only step that does.

\begin{schema}{Quiesce}
  \Xi MailEventState
\end{schema}

%%----------------------------------------------------------------------
\section{Model Checking}
%%----------------------------------------------------------------------

$MailEventInv$ is the state invariant.  Two formulas state delivery:

\begin{quote}
  $AG\,(\,sentMark~k < mark~k \Rightarrow EF\;sentMark~k = mark~k\,)$ for
  each kind, that every arrival is eventually announced; and \\
  $AG\,(\,sentMark~k < mark~k \land wake = \emptyset \Rightarrow
  enabled(PollWake)\,)$, that the announcement does not depend on a push: the
  poll wake needs nothing, so a lost push strands nothing.
\end{quote}

The goal that the properties are not vacuous: an event goes out
($cover~(kMail, 1) = 1$).  Three frozen
variants are the negative controls:

\begin{itemize}
  \item \texttt{notification-mail-event-nolock.tex}: $ReadBegin$ does not take the
    lock.  Two readings take the same mark, both find it newer, and $cover$
    reaches $2$: one arrival, two events.
  \item \texttt{notification-mail-event-markread.tex}: $ReadCommit$ marks the mail
    read when it announces it ($unreadN' = 0$ with no read): the invariant
    $unreadN + readN = arrivedN$ is violated.
  \item \texttt{notification-mail-event-nopoll.tex}: $PollWake$ is never enabled:
    the second delivery formula is \textbf{FALSE}, a lost push strands the
    arrival.
\end{itemize}

%%----------------------------------------------------------------------
%% Reproducing the checks.  Run from the repo root; one probcli at a time.
%%
%%   PROBCLI="${PROBCLI:-$HOME/Applications/ProB/probcli}"
%%   P="-p DEFAULT_SETSIZE 1 -p MAXINT 6 -p MAX_INITIALISATIONS 20 \
%%      -p MAX_OPERATIONS 500 -p TIME_OUT 600000"
%%   ( ulimit -v 8000000; "$PROBCLI" docs/specs/notification-mail-event.tex $P -model_check )
%%   "$PROBCLI" docs/specs/notification-mail-event.tex $P -model_check -ctlformula \
%%       "AG({sentMark(kMail)<mark(kMail)} => EF {sentMark(kMail)=mark(kMail)})"
%%   "$PROBCLI" docs/specs/notification-mail-event.tex $P -model_check -ctlformula \
%%       "AG({sentMark(kMail)<mark(kMail) & card(wake)=0} => e(PollWake))"
%%   goals: "cover(kMail,1)=1"
%%
%% Negative controls:
%%
%%   for V in nolock markread; do
%%     "$PROBCLI" docs/specs/notification-mail-event-$V.tex $P -model_check
%%   done
%%   "$PROBCLI" docs/specs/notification-mail-event-nopoll.tex $P -model_check -ctlformula \
%%       "AG({sentMark(kMail)<mark(kMail) & card(wake)=0} => e(PollWake))"
%%----------------------------------------------------------------------

\end{document}
